\documentclass[12pt]{article}
\usepackage{amsmath, amsthm, amssymb}
\usepackage[utf8]{inputenc}
\usepackage[margin=1in]{geometry}

\newtheorem{theorem}{Theorem}

\begin{document}

\section{Extrema of Finite Sets}
\begin{theorem}
Let $(M, <)$ be the underlying totally ordered set of an o-minimal structure. Every finite, non-empty subset $A \subseteq M$ has a minimum and a maximum.
\end{theorem}
\begin{proof}
We proceed by mathematical induction on the cardinality $n = |A|$.

\textbf{Base case:} For $n = 1$, let $A = \{x\}$. Clearly, $x$ is both the minimum and maximum since $\forall y \in A$, $x \le y$ and $y \le x$.

\textbf{Inductive step:} Assume that every subset of size $k \ge 1$ has a minimum and a maximum. Let $A$ be a subset of size $k + 1$. We write $A = B \cup \{x\}$, where $|B| = k$ and $x \notin B$.

By the inductive hypothesis, $B$ has a maximum $M_B$ and a minimum $m_B$. Since the set is totally ordered, we define:
\begin{align*}
M_A &= \max(M_B, x) \\
m_A &= \min(m_B, x)
\end{align*}

Because $M_A \in A$ and is greater than or equal to all elements in $B$ and equal to or greater than $x$, it is the maximum of $A$. Similarly, $m_A$ is the minimum. By the principle of mathematical induction, this holds for all finite sets.
\end{proof}



\end{document}